\section{Map and Capacity: Fermion Propagation}
\label{sec:map_capacity}

The fermion problem has three parts: construct an exactly norm-preserving update, control its low-momentum dispersion and species content, and identify any mass-generating coupling. The substrate's two $D_4$ sectors motivate the candidate below; they do not alone determine its Hamiltonian, representation or interaction charges.

\subsection{Discrete Update versus Continuum Approximation}
\label{sec:unitary_flux_constraint}
For Hermitian Hamiltonian matrices $\alpha_k,\beta$, a central-difference approximation to the Dirac Hamiltonian is
\begin{equation}
 H\psi=-i\sum_{k=1}^3\alpha_k\frac{\psi(x+e_k\ell)-\psi(x-e_k\ell)}{2\ell}+m\beta\psi.
\end{equation}
Its forward-Euler step is
\begin{equation}
\begin{aligned}
 \psi(t+\tau)={}&\psi(t)-\frac{\tau}{2\ell}\sum_k\alpha_k
 [\psi(x+e_k\ell)-\psi(x-e_k\ell)]\\
 &-i\tau m\beta\psi(t).
\end{aligned}
\label{eq:unitary_flux}
\end{equation}
This is an approximation, not an exactly unitary lattice rule. Even the homogeneous mass-only Euler factor changes the norm at nonzero step (\nolinkurl{DiscreteEvolution}). Formally dividing by $\tau$ gives
\begin{equation}
 i\partial_t\psi=-i\sum_k\alpha_k\partial_k\psi+m\beta\psi.
\label{eq:diff_form}
\end{equation}
The central-difference coefficient $\tau/(2\ell)$ is $1/2$ in link-speed units. It is a kinematic coefficient of this stencil, not a determination of an interaction coupling.

Exactly norm-preserving rotation-and-stream updates exist \cite{meyer_quantum_1996,feynman_quantum_1965}; \nolinkurl{DiscreteEvolution} formalizes a finite-ring example. Its dispersion and twins must be checked separately. Likewise, exponentiating a self-adjoint Hamiltonian gives exact unitary evolution, but this does not prove a finite-range one-tick stencil: an exponential can propagate through repeated local hops. The domain-wall exponential remains conditional on that locality question.

\subsection{The Clifford Constraint}
For a first-order Hamiltonian $H(p,m)=\sum_k p_k\alpha_k+m\beta$, Lean proves the equivalence
\begin{equation}
 H(p,m)^2=(|p|^2+m^2)\mathbb I\quad\text{for all }p,m
\end{equation}
if and only if the generators square to one and pairwise anticommute (\nolinkurl{Dirac.dispersion_iff_clifford}). This is a consequence of the specified dispersion, not of unitarity alone.

For the spacetime convention $\eta=\operatorname{diag}(-,+,+,+)$ one consistent choice is $\gamma^0=\beta$, $\gamma^k=\beta\alpha_k$, with
\begin{equation}
 \{\gamma^\mu,\gamma^\nu\}=-2\eta^{\mu\nu}\mathbb I.
 \label{eq:anti_commute}
\end{equation}
Multiplying the Hamiltonian equation by $\beta$ gives
\begin{equation}
 (i\gamma^\mu\partial_\mu-m)\psi=0.
 \label{eq:covariant_dirac}
\end{equation}
The sign of the Clifford convention must not be mixed with a different metric convention. Deriving the physical representation and its continuum approximation remains distinct from the algebraic equivalence.

\subsection{Twins and a Two-Sector Domain Wall}
The one-coin streaming update obeys $W(k+\pi)=-W(k)$ (\nolinkurl{Doubling.symbol_shift}). Its twin couples identically in every tested link background (\nolinkurl{TwinTests.twin_anticommutes}). A particular two-coin partial-stream update gaps the twin, but its light speed is $|\cos r|$; restoring unit link speed ungaps the twin (\nolinkurl{WalkSpeed}). These results rule out shortcuts within those specified update families.

A different candidate uses the two sectors as two slices of a Wilson/domain-wall Hamiltonian. This construction \emph{does} use a Wilson term; no claim of a Wilson-free mechanism is made. At the chosen point $M=r=1$, its zero-momentum inter-sector hops are $-\beta(1\pm\Gamma)/2$, where $\Gamma=i\beta\alpha_4$. Lean proves that the kernel conditions are
\begin{equation}
 \Gamma x_1=x_1,\qquad \Gamma x_2=-x_2
\end{equation}
(\nolinkurl{DomainWall.zero_modes}). The spatial momentum generators commute with $\Gamma$ and their square gives $|k|^2$ under the stated Clifford hypotheses. This supports the proposed chiral leading response. The full squared operator and necessary energy branches are derived below. A concrete representation is now checked exactly in Lean:
\begin{equation}
 \beta=\begin{pmatrix}0&I_2\\I_2&0\end{pmatrix},\qquad
 \Gamma=\begin{pmatrix}I_2&0\\0&-I_2\end{pmatrix},\qquad
 \alpha_j=\begin{pmatrix}\sigma_j&0\\0&-\sigma_j\end{pmatrix},
\end{equation}
where $\sigma_j$ are the Pauli matrices. The module \nolinkurl{ConcreteFermion} verifies the Clifford relations, chirality commutation and Hermiticity of the full eight-component Hamiltonian for real momenta and detuning. Each chirality projector has rank two. At $k=0$, $\delta=0$, the entire kernel is parametrized by
\begin{equation}
 x_1=(u_0,u_1,0,0),\qquad x_2=(0,0,u_2,u_3),\qquad u\in\mathbb C^4.
\end{equation}
Lean proves this parametrization is a linear equivalence, hence the kernel has complex dimension four: two independent components on each wall, with opposite $\Gamma$ chirality. This is a result about the specified Hamiltonian; persistence has not yet selected this representation or identified these components with physical species.

The script \nolinkurl{calculations/domain_wall_check.py} checks a concrete matrix representation numerically. The corner mass polynomial has a proved positive bound (\nolinkurl{DetuningMass.twin_gap}). The new \nolinkurl{DomainWallSpectrum} result goes beyond corners: for the sine-momentum operator and Wilson detuning
\begin{equation}
 s(k)=\sum_j\sin^2 k_j,\qquad \delta(k)=\sum_j(1-\cos k_j),
\end{equation}
a nonzero zero-energy vector anywhere in $[-\pi,\pi]^3$ requires $k=0$. The proof uses the full matrix polynomial and does not assume the absence of off-corner zeros.

For every real-energy eigenpair it also proves the uniform bound $|E|<8$ in the model's units. The sufficient step condition $0<|\tau|\leq1/2$ then gives $|\tau E|<2\pi$, so a unit eigenphase $e^{-i\tau E}=1$ cannot occur away from $k=0$. This supplies the scalar bandwidth premise on the entire momentum cell. It does not select the physical clock step, construct a finite-stencil exponential update, or assign the calculated zero-mode components to physical chiral gauge representations.

The interpretation of sectors as physical chiral gauge representations remains a separate task. Complementary zero-mode conditions do not prove that one sector is dynamically static or that a particular weak gauge group acts only on it.

\subsection{An Exactly Local Unitary Split Update}
The Hamiltonian calculation can be supplemented by a finite spatial stencil. Let $S_j$ shift one coordinate on a finite periodic cubic lattice and put
\begin{equation}
 A_j=\frac{\beta+i\alpha_j}{2},\qquad
 \widehat A_j=\operatorname{diag}(A_j,A_j).
\end{equation}
Lean checks $A_j^2=0$ and $A_jA_j^\dagger+A_j^\dagger A_j=I$. Consequently
\begin{equation}
 B_j=S_j\otimes\widehat A_j+S_j^\dagger\otimes\widehat A_j^\dagger,
 \qquad B_j^\dagger=B_j,\quad B_j^2=I,
\end{equation}
and $R_j(\tau)=\cos\tau\,I+i\sin\tau\,B_j$ is exactly unitary and reads only the current site and its two neighbors along coordinate $j$ (\nolinkurl{LocalUnitary}, \nolinkurl{SplitUpdate}). With $H_0=H(0,3)$, define
\begin{equation}
 U_\tau=(I_{\rm space}\otimes e^{-i\tau H_0})R_0(\tau)R_1(\tau)R_2(\tau).
\end{equation}
The on-site exponential mixes only eight internal components. Lean proves unitarity and conservation of the total squared amplitude on every finite periodic cubic lattice, together with a precise input-dependence theorem: one tick reads at most $3^3=27$ spatial sites, independently of lattice size (\nolinkurl{PropagationBounds}, \nolinkurl{WilsonSplit}).

The directional phase symbol is
\begin{equation}
 B_j(k_j)=\operatorname{diag}(\cos k_j\,\beta-\sin k_j\,\alpha_j,
                              \cos k_j\,\beta-\sin k_j\,\alpha_j).
\end{equation}
The split generators satisfy $H_0-\sum_j B_j(k_j)=H(k,\delta(k))$. The update's momentum symbol satisfies $U_0(k)=I$ and has first derivative $-iH(k,\delta(k))$ at $\tau=0$ (\nolinkurl{SplitGenerator}). This is a first-order consistency result. It does not establish equality with $e^{-i\tau H(k)}$ at nonzero tick size, a many-tick convergence bound, or the finite-tick zero-mode and eigenphase multiplicities. In particular the Hamiltonian bandwidth argument above cannot simply be transferred to this product. A separate double-precision diagnostic at $k=0$, $\tau=0.1$ reports a nearest eigenphase magnitude of approximately $7.47\times10^{-4}$ rather than zero. This is a numerical warning, not a Lean obstruction theorem: the displayed ordering should not be treated as an exact vacuum-preserving physical update. This motivates separating the on-site factors as follows.

\paragraph{Vacuum-preserving ordering.}
Write $W=H(0,0)$ for the wall hopping alone, $\beta_8=\operatorname{diag}(\beta,\beta)$ and
\begin{equation}
 R_\beta(\tau)=\cos\tau\,I_8+i\sin\tau\,\beta_8.
\end{equation}
The corrected on-site coin and update are
\begin{equation}
 \begin{aligned}
 \widetilde Q_\tau&=e^{-i\tau W}[R_\beta(\tau)^\dagger]^3,\\
 \widetilde U_\tau&=(I_{\rm space}\otimes\widetilde Q_\tau)
                      R_0(\tau)R_1(\tau)R_2(\tau).
 \end{aligned}
\end{equation}
At $k=0$, each spatial factor has symbol $R_\beta(\tau)$, so Lean proves the exact identity
\begin{equation}
 \widetilde U_\tau(0)=e^{-i\tau W}
 \qquad\hbox{for every real }\tau.
\end{equation}
Since $W$ annihilates the previously characterized four-dimensional vacuum kernel, every vector in that kernel is fixed exactly by $\widetilde U_\tau(0)$ (\nolinkurl{VacuumSplit.vacuum_preserved}). The proof uses the vacuum projector and the exponential, with no small-tick approximation. The module \nolinkurl{VacuumUniform} also proves this for the corresponding spatially constant states on every finite periodic lattice, connecting the symbol identity to the actual lattice update.

The corrected update is exactly unitary, conserves total squared amplitude, and has the same at-most-27-site spatial input bound. Separating the on-site factors changes only terms beyond first order: Lean checks $\widetilde U_0(k)=I$ and $\partial_\tau\widetilde U_\tau(k)|_{\tau=0}=-iH(k,\delta(k))$ (\nolinkurl{VacuumGenerator}). Thus exact preservation of the known vacuum modes, unitarity, finite propagation and the stated Wilson tangent hold simultaneously in this construction.

This proves preservation of the known four modes; it does not prove they exhaust the fixed space of the finite-tick update. Additional unit eigenphases, including those away from the origin, a permissible physical tick, many-tick error control and physical species identification remain separate obligations. These constraints also do not establish that persistence uniquely selects this ordering.

\subsection{Mass and the Scalar Bridge}
\label{sec:capacity_mechanism}
Under the stated Clifford relations the hop products reduce to complementary projectors (\nolinkurl{DetuningMass.hop_products}). Writing the diagonal operator as $K+\delta\beta$, Lean now proves the full squared matrix has diagonal blocks $(s+\delta^2)I+P_-$ and $(s+\delta^2)I+P_+$ and off-diagonal blocks $-\delta I$, with all momentum-hop cross terms cancelling (\nolinkurl{DomainWallSpectrum.hamiltonian_square}). On the common $\Gamma=+1$ subspace this gives the reduced block
\begin{equation}
 \begin{pmatrix}\delta^2&-\delta\\-\delta&\delta^2+1\end{pmatrix}.
\end{equation}
Lean checks that
\begin{equation}
 \begin{aligned}
 m(\delta)&=\sqrt{\delta^2+1/4}-1/2,\\
 m(m+1)&=\delta^2,\qquad 0\leq m\leq\delta^2.
 \end{aligned}
\end{equation}
The full operator also obeys a polynomial implying, for every nonzero real-energy eigenvector,
\begin{equation}
 E^2=s+m(\delta)^2\quad\hbox{or}\quad E^2=s+(m(\delta)+1)^2
\end{equation}
(\nolinkurl{DomainWallSpectrum.energy_branches}). These are necessary spectral branches derived from the Hamiltonian, rather than roots postulated independently. For the concrete representation above, \nolinkurl{ConcreteSpectrum.square_charpoly} strengthens this necessary condition to the exact identity
\begin{equation}
 \det(\lambda I_8-H^2)=
 [\lambda-s-m(\delta)^2]^4
 [\lambda-s-(m(\delta)+1)^2]^4.
\end{equation}
The two roots are strictly separated, since $(m+1)^2-m^2=2m+1>0$. Lean proves each has algebraic multiplicity four (\nolinkurl{ConcreteSpectrum.square_branch_multiplicities}). Thus both squared-energy branches occur in this model. The signed count is now established as well. A real matrix $J$ with $J^2=-I$ obeys $JH^T=-HJ$, giving equal characteristic polynomials for $H$ and $-H$. Combining this pairing with the squared polynomial yields
\begin{equation}
 \det(E I_8-H)=[E^2-s-m^2]^2[E^2-s-(m+1)^2]^2
\end{equation}
(\nolinkurl{SignedSpectrum.signed_charpoly}). Whenever the light energy is nonzero, each of $\pm\sqrt{s+m^2}$ has algebraic multiplicity two; each of $\pm\sqrt{s+(m+1)^2}$ always has multiplicity two. At the origin the polynomial is $E^4(E-1)^2(E+1)^2$, consistent with the explicit four-dimensional zero kernel (\nolinkurl{EnergyCounts}). These counts do not identify physical species. The detuning of each physical fermion is not yet derived; on-site detuning must not be identified with a change of the hop without an additional calculation.

A phase-covariant scalar link is a candidate interaction between sector states. Given representations and charges, its continuum operator can take the Yukawa form
\begin{equation}
 \mathcal L_{\rm CAP}=-y(\bar\psi_L\phi\psi_R+\bar\psi_R\phi^\dagger\psi_L).
 \label{eq:yukawa}
\end{equation}
The charge-pattern checks in Paper III are conditional on the chosen representations and scalar link. Identifying a background scalar with a linear physical mass $y\langle\phi\rangle$, and reconciling that law with the detuning spectrum, require a derived field and scale map. Neither a numerical Yukawa value nor its normalization follows from the hopping stencil alone.

The magnetic-moment response also remains to be calculated from the update; \cref{app:impedance_power_transfer} records the failure of the earlier winding-angle argument.


\paragraph{Comparison cost versus the fermion kinetic operator.}
\nolinkurl{CommonResponse.wilson_form} rewrites the positive comparison penalty as
$\operatorname{Re}\langle\psi,(2-U-U^\dagger)\psi\rangle/4$.
This is a transport-mismatch quadratic form. It does not alone derive the first-order
Dirac action, fermionic statistics or the identification of glue-class labels with the
spinor components used above. Those remain explicit carrier-to-response obligations.
